\documentclass[11pt]{amsart}
\usepackage[T1]{fontenc}
\usepackage{lmodern,microtype,amsmath,amssymb,amsthm,mathtools,enumitem}
\usepackage[margin=1.1in]{geometry}
\usepackage[hidelinks]{hyperref}
\usepackage{xurl}

\newtheorem{theorem}{Theorem}[section]
\newtheorem{lemma}[theorem]{Lemma}
\newtheorem{corollary}[theorem]{Corollary}
\newtheorem{proposition}[theorem]{Proposition}
\theoremstyle{definition}
\newtheorem{definition}[theorem]{Definition}
\theoremstyle{remark}
\newtheorem{remark}[theorem]{Remark}
\numberwithin{equation}{section}

\newcommand{\PP}{\mathbb P}
\newcommand{\E}{\mathbb E}
\newcommand{\CM}{\mathrm{CM}}
\newcommand{\eps}{\varepsilon}
\newcommand{\cF}{\mathcal F}
\newcommand{\cC}{\mathcal C}
\newcommand{\dmax}{d_{\max}}

\hypersetup{pdftitle={Hub paths and the largest component of the barely subcritical configuration model},pdfauthor={Yitzchak Shmalo}}

\title[Hub paths in the barely subcritical configuration model]{Hub paths and the largest component\\ of the barely subcritical configuration model}
\author{Yitzchak Shmalo}
\address{Einstein Institute of Mathematics, The Hebrew University of Jerusalem, Jerusalem, Israel}
\email{yitzchak.shmalo@gmail.com}
\date{18 October 2026}
\subjclass[2020]{05C80, 60C05}
\keywords{configuration model, random graph with given degrees, barely subcritical, largest component, high-degree vertices}

\begin{document}
\begin{abstract}
For the configuration model with a degree sequence whose third moment converges, write $\eps_n=1-\nu_n$, where $\nu_n\to1$ from below is the size-biased mean minus one. Van der Hofstad asked to show that in the barely subcritical regime the largest component has order $\max\{\dmax/\eps_n,\ \eps_n^{-2}\log(n\eps_n^3)\}$, and predicted that it has size $(\dmax/\eps_n)(1+o_\PP(1))$ when $\eps_n$ tends to zero sufficiently slowly. We show that neither holds in this generality. We use degree sequences with vertices of degree one and three and a polynomial number of vertices of a large common degree $d$. Paths of $K$ such hubs joined by direct edges then appear, and the component of such a path is a tree with about $Kd/\eps_n$ vertices. Taking $K=2$ gives $L_1\ge(2-\eta)\dmax/\eps_n$ with high probability, for $\eps_n$ tending to zero as slowly as we like; letting $K\to\infty$ gives $L_1/\max\{\dmax/\eps_n,\eps_n^{-2}\log(n\eps_n^3)\}\to\infty$ in probability. In all examples the degree of a uniform vertex converges in law to a variable on $\{1,3\}$ with all three moments, and $n\eps_n^3\to\infty$.
\end{abstract}
\maketitle

\section{Introduction}

Let $\mathbf d=(d_1,\dots,d_n)$ be a sequence of positive integers with even sum $\ell=\sum_i d_i$. The configuration model $\CM_n(\mathbf d)$ attaches $d_i$ half-edges to vertex $i$ and pairs the $\ell$ half-edges by a uniformly random perfect matching. Let $D_n$ be the degree of a uniformly chosen vertex and
\[
\nu_n=\frac{\E[D_n(D_n-1)]}{\E[D_n]}=\frac1\ell\sum_{i}d_i(d_i-1),\qquad \eps_n=1-\nu_n=-\frac1\ell\sum_i d_i(d_i-2).
\]
We write $L_1$ for the number of vertices of the largest component of $\CM_n(\mathbf d)$ and $\dmax=\max_i d_i$. The barely subcritical regime is $\eps_n\to0$ with $n\eps_n^3\to\infty$, that is, $\nu_n$ tends to one from below but stays outside the critical window.

In his Saint-Flour lecture notes~\cite{vdH17}, van der Hofstad recalls Janson's theorem~\cite{Jan08} that for power-law degrees in the strictly subcritical regime the largest component has size $\dmax/(1-\nu)(1+o_\PP(1))$, and writes that one could expect this behaviour to persist when $\nu_n\to1$ sufficiently slowly and $\dmax$ is sufficiently large; ``we predict that the maximal cluster has size $\dmax/\eps_n(1+o_\PP(1))$'' when $\eps_n$ tends to zero sufficiently slowly. Close to the critical window one expects instead the Erd\H{o}s--R\'enyi order $\eps_n^{-2}\log(n\eps_n^3)$ (see \cite{Luc90} for the random graph). This led to~\cite[Open Problem~4.13]{vdH17}: under convergence in distribution of $D_n$ to a limit $D$ and convergence of the first two moments \cite[Condition~1.6(a)--(c)]{vdH17}, and $\E[D_n^3]\to\E[D^3]$, identify the size of the largest component in the barely subcritical regime and
\begin{quote}
show that the largest connected component has a size of order
\[
\max\{\dmax/\eps_n,\ \eps_n^{-2}\log(n\eps_n^3)\},
\]
thus confirming that this setting has a phase transition.
\end{quote}
The notes use $\eps_n=\nu_n-1<0$; we use the positive quantity $1-\nu_n$, which is what the displayed order requires. The notes to the corresponding chapter of the book draft~\cite{vdH26} contain the same prediction $L_1=(\dmax/\eps_n)(1+o_\PP(1))$ but not the open problem.

We show that both statements fail. Write
\[
H_n=\frac{\dmax}{\eps_n},\qquad E_n=\eps_n^{-2}\log(n\eps_n^3).
\]

\begin{theorem}\label{thm:examples}
\begin{enumerate}[label=\textup{(\alph*)}]
\item Let $\varphi$ be any function with $\varphi(n)\to0$. There are degree sequences $\mathbf d$, along a sequence $n\to\infty$, such that $D_n$ converges in distribution to $D$ with $\PP(D=1)=3/4$ and $\PP(D=3)=1/4$, $\E[D_n^j]\to\E[D^j]$ for $j=1,2,3$, $\eps_n\to0$, $\eps_n\ge\varphi(n)$, $n\eps_n^3\to\infty$, $H_n/E_n\to\infty$, and for every $\eta>0$
\[
\PP\bigl(L_1\ge(2-\eta)\,\dmax/\eps_n\bigr)\to1 .
\]
In particular $L_1\ne(\dmax/\eps_n)(1+o_\PP(1))$.
\item There are degree sequences with the same limit law and the same three convergent moments, $\eps_n\to0$ and $n\eps_n^3\to\infty$, such that
\[
\frac{L_1}{\max\{H_n,E_n\}}\to\infty\quad\text{in probability};
\]
in fact $L_1\ge c\sqrt{\log n}\,\max\{H_n,E_n\}$ with probability tending to one, for an absolute constant $c>0$.
\end{enumerate}
\end{theorem}

Both parts are instances of Theorem~\ref{thm:main} below, which applies to degree sequences with vertices of degree $1$ and $3$ and $M$ vertices of a common degree $d$, the hubs. Its mechanism is simple. When $Md^2/\ell$ is small compared with $\eps_n$, a hub sees another hub through a direct edge with probability about $Md^2/\ell$, while the component it would have on its own is a subcritical tree of size about $d/\eps_n$. Paths of $K$ hubs joined by direct edges exist as soon as $M(Md^2/\ell)^{K-1}\to\infty$, and the component of such a path is, with high probability, a tree that contains no further hub and has at least $(1-\eta)Kd/\eps_n$ vertices. A second-moment argument over these paths finishes the proof. The size estimate needs some care, since the exploration of the component starts with $Kd$ active half-edges and runs for about $Kd/\eps_n$ steps; we use a compensated process whose increments stay bounded at hub discoveries. The hubs contribute $Md^3/n\to0$ to the third moment, so they are invisible in the limit law.

The same mechanism appears in the work of Coulson and Perarnau~\cite{CP23}. For a strictly subcritical sequence with about $n/\Delta^2$ vertices of degree $\Delta$ and the rest of degree one, they compare the graph induced on the high-degree vertices with an Erd\H{o}s--R\'enyi graph and obtain a lower bound for $L_1$ that matches their upper bounds \cite[Proposition~1.12 and Remark~1.13]{CP23}. In their examples $R=\frac1\ell\sum_i d_i(d_i-2)^2\sim\Delta\to\infty$, so the third moment is unbounded and Open Problem~4.13 does not apply. Here the third moment converges, the graph induced on the hubs is much sparser, and we use paths of a bounded or slowly growing number of hubs joined by direct edges. We do not determine the order of $L_1$ in our examples, and we do not propose a replacement for the order in Open Problem~4.13; Theorem~\ref{thm:main} only shows that the number of hubs a component can collect enters the answer.

\section{The degree sequences and the main theorem}

All asymptotic statements refer to $n\to\infty$ along the sequences considered. For integers $n,M,d,y\ge1$ with $4\mid n$ and $4\mid(d-1)^2M$ put
\begin{equation}\label{eq:counts}
N_3=\frac n4-\frac{(d-1)^2M}{4}-y,\qquad N_1=n-N_3-M,
\end{equation}
and assume $N_1,N_3\ge0$ and $d\ge4$. Let $\mathbf d=\mathbf d(n,M,d,y)$ have $N_1$ vertices of degree one, $N_3$ vertices of degree three, and $M$ vertices of degree $d$, called hubs.

\begin{lemma}\label{lem:counts}
The total degree $\ell=N_1+3N_3+dM$ is even, and
\[
\ell\eps_n=4y,\qquad \ell=\frac{3n}{2}-\frac{(d-1)^2M}{2}+(d-1)M-2y .
\]
\end{lemma}

\begin{proof}
We have $\ell=n+2N_3+(d-1)M$ and $(d-1)M$ is even, since $4\mid(d-1)^2M$. Next, $\sum_i d_i(d_i-2)=-N_1+3N_3+d(d-2)M=-n+4N_3+(d-1)^2M=-4y$, which is $-\ell\eps_n$. The formula for $\ell$ follows by substituting~\eqref{eq:counts}.
\end{proof}

Let $K\ge2$ be an integer, possibly depending on $n$, put $q=K-1$, and define
\begin{equation}\label{eq:mu}
P_r=\prod_{j=0}^{r-1}(\ell-2j-1),\qquad w=d^2\bigl(d(d-1)\bigr)^{K-2},\qquad \mu=\frac{(M)_K\,w}{P_q},
\end{equation}
where $(M)_K=M(M-1)\cdots(M-K+1)$. We write $\eps=\eps_n$ from now on.

\begin{theorem}\label{thm:main}
Let $\mathbf d=\mathbf d(n,M,d,y)$ and $K\ge2$ be as above, all depending on $n$, and assume
\begin{enumerate}[label=\textup{(C\arabic*)}]
\item\label{C1} $y=o(n)$ and $Md^3=o(n)$;
\item\label{C2} $n\eps^3\to\infty$ and $d\eps/\log(n\eps^3)\to\infty$;
\item\label{C3} $K^2d^2+KMd^2=o(n\eps)$;
\item\label{C4} $\mu\to\infty$.
\end{enumerate}
Then $D_n$ converges in distribution to $D$ with $\PP(D=1)=3/4$, $\PP(D=3)=1/4$, and $\E[D_n^j]\to\E[D^j]$ for $j=1,2,3$; moreover $\eps\to0$, $\dmax=d$ and $H_n/E_n\to\infty$. For every fixed $\eta\in(0,1)$,
\[
\PP\Bigl(L_1\ge(1-\eta)\,\frac{K\dmax}{\eps}\Bigr)\to1 .
\]
\end{theorem}

The limit law has $\E[D]=3/2$, $\E[D^2]=3$, $\E[D^3]=15/2$ and $\E[D(D-1)]/\E[D]=1$, and $\PP(D=2)=0$, so it is a critical law with finite third moment and no vertices of degree two.

\begin{proof}[Proof of the first assertions]
By~\ref{C1}, $(d-1)^2M\le Md^3=o(n)$ and $y=o(n)$, so $N_3/n\to1/4$, $M/n\to0$ and $N_1/n\to3/4$. For $j\le3$, $\E[D_n^j]=(N_1+3^jN_3+d^jM)/n$ and $d^jM\le d^3M=o(n)$, which gives the convergence of the law and of the three moments. By Lemma~\ref{lem:counts}, $\ell=(3/2+o(1))n$ and $\eps=4y/\ell\to0$. Finally $H_n/E_n=d\eps/\log(n\eps^3)\to\infty$ by~\ref{C2}.
\end{proof}

We collect the consequences of the hypotheses used below. Since $\ell\sim3n/2$, \ref{C3} holds with $\ell$ in place of $n$. By~\ref{C2}, $d\eps\to\infty$; in particular $d\to\infty$. Hence
\begin{equation}\label{eq:cons}
\frac{Kd}{\ell\eps^2}=\frac{K^2d^2}{\ell\eps}\cdot\frac1{Kd\eps}\to0,\qquad \frac{Md^2}{\ell\eps}\to0,\qquad \frac{K^2}{\ell}\to0,\qquad \ell\eps\to\infty,
\end{equation}
the third because $K^2\le K^2d^2=o(\ell\eps)$ and $\eps<1$, the last because $\ell\eps\ge n\eps^3$ (note $\ell\ge n$ and $\eps<1$).

\section{Hub paths and their moments}

\begin{definition}
A \emph{witness} is a sequence $(v_1,\dots,v_K)$ of distinct hubs together with half-edges $a_i$ at $v_i$ and $b_{i+1}$ at $v_{i+1}$ for $1\le i<K$, such that $a_i\ne b_i$ for $1<i<K$. It is \emph{present} if $a_i$ is paired with $b_{i+1}$ for every $i<K$. Let $Z$ be the number of present witnesses and let $D_0$ be the number of ordered pairs of present witnesses whose hub sets are disjoint.
\end{definition}

For a given sequence of hubs there are $d$ choices for $a_1$, $d$ for $b_K$ and $d(d-1)$ for each interior pair $(b_i,a_i)$, that is, $w$ choices in all.

\begin{lemma}\label{lem:moments}
$\E Z=\mu$ and $\E D_0\le(1+o(1))\mu^2$.
\end{lemma}

\begin{proof}
A prescribed set of $r$ disjoint pairs of half-edges is contained in the uniform matching with probability $1/P_r$: given the first $j$ prescribed pairs, the next one occurs with probability $1/(\ell-2j-1)$. A witness prescribes $q$ disjoint pairs, and there are $(M)_Kw$ witnesses, so $\E Z=\mu$. Two witnesses with disjoint hub sets use disjoint half-edges and prescribe $2q$ disjoint pairs, and there are $(M)_{2K}w^2$ ordered pairs of them. Hence
\[
\frac{\E D_0}{\mu^2}=\frac{(M)_{2K}}{(M)_K^2}\cdot\frac{P_q^2}{P_{2q}}\le\prod_{j=0}^{q-1}\frac{\ell-2j-1}{\ell-2q-2j-1}\le\Bigl(1+\frac{2q}{\ell-4q}\Bigr)^{q}\le\exp\Bigl(\frac{2q^2}{\ell-4q}\Bigr),
\]
using $(M)_{2K}=(M)_K(M-K)_K\le(M)_K^2$. The right side tends to one by~\eqref{eq:cons}.
\end{proof}

\section{The component of a present witness}

Fix a witness $W$ and condition on the event that it is present. The remaining $\ell-2q$ half-edges are then matched uniformly at random. Let $\cC(W)$ be the component containing the hubs of $W$. All bounds below are uniform in $W$.

\smallskip\noindent\emph{The exploration.} At time $0$ the $K$ hubs of $W$ are discovered, their $A_0=Kd-2q$ unpaired half-edges are active, and every other half-edge is sleeping. At time $t$, if $A_t\ge1$, we take the active half-edge with the smallest label and reveal its partner, which is uniform among the
\[
B_t=\ell-2q-2t-1
\]
other unpaired half-edges. If the partner is active, the two are paired and $A_{t+1}=A_t-2$; we call this a \emph{surplus} step. If the partner belongs to a sleeping vertex $v$, then $v$ is discovered, the pair is formed, the other $d_v-1$ half-edges of $v$ become active, and $A_{t+1}=A_t+d_v-2$. Let $\tau=\min\{t:A_t=0\}$. The vertices discovered by time $\tau$ form $\cC(W)$. Let $\cF_t$ be the $\sigma$-field of the first $t$ steps, let $R_t$ be the set of vertices discovered during the first $t$ steps (not including the hubs of $W$), and let $\Delta A_t=A_{t+1}-A_t$.

\begin{lemma}\label{lem:drift}
On $\{A_t\ge1\}$,
\begin{align}
\E[\Delta A_t\mid\cF_t]&=\frac{-\ell\eps-Kd(d-2)-\sum_{v\in R_t}d_v(d_v-2)-2(A_t-1)}{B_t},\label{eq:driftA}\\
\E[(\Delta A_t)^2\mid\cF_t]&=\frac{\sum_{v\ \mathrm{sleeping}}d_v(d_v-2)^2+4(A_t-1)}{B_t}.\label{eq:varA}
\end{align}
\end{lemma}

\begin{proof}
A sleeping vertex $v$ offers $d_v$ half-edges, each giving $\Delta A_t=d_v-2$, and the other $A_t-1$ active half-edges each give $-2$. The sum of $d_v(d_v-2)$ over sleeping vertices equals $\sum_id_i(d_i-2)=-\ell\eps$ minus the contributions of the $K$ hubs of $W$ and of $R_t$.
\end{proof}

Let $T_0=\lfloor\ell\eps/4\rfloor$, let $\sigma$ be the first time $t$ at which $R_t$ contains a hub ($\sigma=\infty$ if there is none), and let $S=\min\{\tau,\sigma,T_0\}$. Since $\eps\to0$ and $q=o(\ell)$, we have $B_t\ge\ell/2$ for $t\le T_0$ and $n$ large.

For $t<S$ the set $R_t$ consists of at most $t$ vertices of degree one or three, so $-\sum_{v\in R_t}d_v(d_v-2)\le t$ (as $k(k-2)\ge-1$), and $A_t\ge1$. The numerator in~\eqref{eq:driftA} is then at most $-\ell\eps+T_0\le-\tfrac34\ell\eps$, and as $B_t\le\ell$,
\begin{equation}\label{eq:betaneg}
\beta_t:=\E[\Delta A_t\mid\cF_t]\le-\eps/2\qquad(t<S).
\end{equation}
In~\eqref{eq:varA}, $k(k-2)^2\le k^3$ gives $\sum_{v}d_v(d_v-2)^2\le N_1+27N_3+Md^3\le28n$ for large $n$ by~\ref{C1}, and $A_t\le\ell$, $n\le\ell$, $B_t\ge\ell/2$; hence
\begin{equation}\label{eq:varbound}
\E[(\Delta A_t)^2\mid\cF_t]\le64\qquad(t<S).
\end{equation}
This holds also when the step at time $t$ discovers a hub, since the bound only uses the state at time $t$.

\begin{lemma}\label{lem:good}
Let $G(W)$ be the event that $\tau<T_0$ and that the exploration discovers no hub and makes no surplus step before $\tau$. Then $\PP(G(W)\mid W\text{ present})\to1$. On $G(W)$ the component $\cC(W)$ is a tree containing exactly $K$ hubs and $|\cC(W)|=K+\tau$.
\end{lemma}

\begin{proof}
All probabilities are conditional on $W$ being present. The time $S$ is bounded by $T_0$, so $A_{t\wedge S}-\sum_{s<t\wedge S}\beta_s$ is a martingale and optional stopping at $T_0$ is elementary. Since $A_S\ge0$, \eqref{eq:betaneg} gives
\begin{equation}\label{eq:ES}
\E S\le\frac{2A_0}{\eps}\le\frac{2Kd}{\eps}.
\end{equation}
For $t<S$, $\E[A_{t+1}^2-A_t^2\mid\cF_t]=2A_t\beta_t+\E[(\Delta A_t)^2\mid\cF_t]\le-\eps A_t+64$; summing up to $S$ and using $A_S^2\ge0$,
\begin{equation}\label{eq:area}
\eps\,\E\sum_{t<S}A_t\le A_0^2+64\,\E S\le K^2d^2+\frac{128Kd}{\eps}.
\end{equation}
At a time $t<S$ the partner is a hub half-edge with conditional probability at most $Md/B_t\le2Md/\ell$, and it is active with conditional probability $(A_t-1)/B_t\le2A_t/\ell$. Since $\{t<S\}\in\cF_t$, summing these conditional probabilities over $t<S$ bounds the probability that some step before $S$ discovers a hub by
\[
\frac{2Md}{\ell}\,\E S\le\frac{4KMd^2}{\ell\eps}\to0,
\]
and the probability of a surplus step before $S$ by
\[
\frac2\ell\,\E\sum_{t<S}A_t\le\frac{2K^2d^2}{\ell\eps}+\frac{256Kd}{\ell\eps^2}\to0,
\]
both by~\ref{C3} and~\eqref{eq:cons}. Finally $\PP(S=T_0)\le\E S/T_0\le(2Kd/\eps)/(\ell\eps/4-1)\to0$ by~\eqref{eq:cons}. Outside these three events no hub is discovered before $S$, so $S\ne\sigma$; also $S\ne T_0$; hence $S=\tau<T_0$, and no step before $\tau$ discovers a hub or is a surplus step. Then each of the $\tau$ steps adds one new vertex and one edge to the path of $W$, which has $K$ vertices and $K-1$ edges, so $\cC(W)$ is a tree with $K+\tau$ vertices whose hubs are those of $W$.
\end{proof}

\begin{lemma}\label{lem:size}
For every fixed $\eta\in(0,1)$,
\[
\PP\Bigl(G(W)\ \text{and}\ |\cC(W)|<(1-\eta)\frac{Kd}{\eps}\Bigm| W\text{ present}\Bigr)\to0 .
\]
\end{lemma}

\begin{proof}
Let $H_t$ be the number of hubs discovered by time $t$, including the $K$ hubs of $W$, and $C_t=A_t-dH_t$, so $C_0=-2q$. A step that discovers a vertex of degree one, a vertex of degree three, or a hub changes $C$ by $-1$, $+1$, $-2$ respectively, and a surplus step changes it by $-2$. Let $r_1(t)$ and $r_3(t)$ count the vertices of degree one and three in $R_t$. The sleeping half-edges are $N_1-r_1(t)$ at vertices of degree one, $3(N_3-r_3(t))$ at vertices of degree three and $d(M-H_t)$ at hubs, so on $\{A_t\ge1\}$
\begin{equation}\label{eq:driftC}
\begin{split}
\gamma_t:=\E[\Delta C_t\mid\cF_t]&=\frac{-(N_1-r_1)+3(N_3-r_3)-2d(M-H_t)-2(A_t-1)}{B_t}\\
&=\frac{-\ell\eps-Md^2+r_1-3r_3-2C_t+2}{B_t},
\end{split}
\end{equation}
where we used $A_t=C_t+dH_t$ and $-N_1+3N_3+d(d-2)M=-\ell\eps$; the terms in $H_t$ cancel.

Put $T=\lfloor(1-\eta/2)Kd/\eps\rfloor$ and $U=\tau\wedge T$. For $t<U$ we have $r_3(t)\le t$ and $C_t\le C_0+t\le t$, since $C$ increases only by one at a time. The numerator in~\eqref{eq:driftC} is therefore at least $-\ell\eps-Md^2-5T$, and $B_t\ge\ell-2q-2T-1$. Hence, whether the numerator is negative or not,
\[
\gamma_t\ge-(1+\rho)\eps\quad(t<U),\qquad 1+\rho=\frac{1+(Md^2+5T)/(\ell\eps)}{1-(2q+2T+1)/\ell},
\]
and $\rho\to0$ by~\eqref{eq:cons}, since $T\le Kd/\eps$. The process $J_t=C_{t\wedge U}-C_0-\sum_{s<t\wedge U}\gamma_s$ is a martingale with $J_0=0$ and increments bounded by $4$ in absolute value. On $\{\tau\le T\}$ we have $C_\tau=-dH_\tau\le-Kd$, so
\[
J_T\le-Kd+2q+(1+\rho)\eps\tau\le-Kd+2K+(1+\rho)(1-\eta/2)Kd\le-\frac{\eta Kd}{4}
\]
for $n$ large, because $\rho\to0$ and $d\to\infty$. By the Azuma--Hoeffding inequality~\cite{Azu67},
\[
\PP(\tau\le T)\le\PP\Bigl(J_T\le-\frac{\eta Kd}4\Bigr)\le\exp\Bigl(-\frac{(\eta Kd/4)^2}{32\,T}\Bigr)\le\exp\Bigl(-\frac{\eta^2}{512}Kd\eps\Bigr)\to0 .
\]
On $G(W)$, $|\cC(W)|=K+\tau$ by Lemma~\ref{lem:good}, and $|\cC(W)|<(1-\eta)Kd/\eps$ forces $\tau<(1-\eta)Kd/\eps\le T$ for $n$ large, since $Kd/\eps\to\infty$.
\end{proof}

\section{Proof of Theorem~\ref{thm:main}}

Fix $\eta\in(0,1)$. Call a witness \emph{excellent} if it is present, $G(W)$ holds and $|\cC(W)|\ge(1-\eta)Kd/\eps$, and let $Y$ be the number of excellent witnesses. By Lemmas~\ref{lem:moments}, \ref{lem:good} and~\ref{lem:size}, uniformly in $W$,
\[
\E Y=\sum_W\PP(W\text{ present})\,\PP(W\text{ excellent}\mid W\text{ present})=(1-o(1))\mu .
\]

\begin{lemma}\label{lem:pairs}
Pointwise, $Y^2\le2Y+D_0$.
\end{lemma}

\begin{proof}
$Y^2$ counts ordered pairs $(W,W')$ of excellent witnesses. If $W$ and $W'$ lie in different components their hub sets are disjoint, and the pair is counted in $D_0$. Suppose they lie in the same component $\cC$. Then $\cC$ is a tree with exactly $K$ hubs, and $W'$ visits $K$ distinct hubs of $\cC$, so its hub set is the hub set of $W$. The path of $W$ is a connected subgraph of $\cC$ whose vertices are these $K$ hubs, and in a tree the smallest connected subgraph containing a given vertex set is unique; so the paths of $W$ and $W'$ have the same vertices and, being contained in a tree, the same edges. A tree has no multiple edges, so the half-edges of $W'$ are determined by its order of the hubs along that path, which is one of the two orientations. Hence for each excellent $W$ at most two excellent $W'$ share its component, and $Y^2\le2Y+D_0$.
\end{proof}

By Lemmas~\ref{lem:pairs} and~\ref{lem:moments} and~\ref{C4},
\[
\frac{\E Y^2}{(\E Y)^2}\le\frac{2\E Y+\E D_0}{(\E Y)^2}\le\frac{2}{(1-o(1))\mu}+\frac{(1+o(1))\mu^2}{(1-o(1))^2\mu^2}\to1,
\]
so $\PP(Y>0)\ge(\E Y)^2/\E Y^2\to1$ by the Cauchy--Schwarz inequality. On $\{Y>0\}$ there is a component with at least $(1-\eta)Kd/\eps$ vertices. This proves Theorem~\ref{thm:main}. \qed

\begin{remark}
The argument uses only first and second moments of witness counts and an exploration started from a prescribed path; no independence between components or between hubs is assumed. The lower bound counts only components with exactly $K$ hubs. Components that collect more hubs, or hubs connected through vertices of degree three, can only make $L_1$ larger.
\end{remark}

\section{The two families}

\begin{corollary}\label{cor:paths}
Let $m$ run through the even integers and let $K=K_m\ge2$ be integers with $12^{K_m}/m\to0$. Put
\[
d=m^2,\qquad M=m^{3K-2},\qquad n=8m^{3K+5},\qquad y=\frac{3n}{8m}.
\]
Then the hypotheses of Theorem~\ref{thm:main} hold, $\eps=(1+o(1))/m$, $\mu=(1+o(1))\,m/12^{K-1}$, and $\PP(L_1\ge(1-\eta)KH_n)\to1$ for every $\eta\in(0,1)$.
\end{corollary}

\begin{proof}
Here $4\mid n$, and $16\mid M$ since $m$ is even and $3K-2\ge4$. We have $y/n=3/(8m)$, $(d-1)^2M/n\le1/(8m^3)$ and $(d-1)M/n\le1/(8m^5)$, so $N_3,N_1>0$ for large $m$, and by Lemma~\ref{lem:counts} $\ell=\tfrac{3n}{2}(1+O(1/m))$ and $\ell\eps=4y=3n/(2m)$, so $\eps=(1+O(1/m))/m$.

\ref{C1}: $y/n\to0$ and $Md^3/n=1/(8m)\to0$. \ref{C2}: $n\eps^3=(8+o(1))m^{3K+2}\to\infty$; since $K\le\log m/\log12$, $\log(n\eps^3)=O(\log^2m)$, while $d\eps=(1+o(1))m$. \ref{C3}: $K^2d^2/(n\eps)=(1+o(1))K^2/(8m^{3K})\to0$ and $KMd^2/(n\eps)=(1+o(1))K/(8m^2)\to0$. \ref{C4}: we have $(M)_K=M^K(1+O(K^2/M))$, $w=d^{2K-2}(1-1/d)^{K-2}=d^{2K-2}(1+O(K/m^2))$ and $P_q=\ell^q(1+O(K^2/\ell))$; also $\ell=12m^{3K+5}(1+O(1/m))$, so $\ell^q=(12m^{3K+5})^{K-1}(1+O(K/m))$. As $K=O(\log m)$, all error factors tend to one and
\[
\mu=(1+o(1))\frac{m^{K(3K-2)}\,m^{4(K-1)}}{12^{K-1}m^{(3K+5)(K-1)}}=(1+o(1))\frac{m}{12^{K-1}}\to\infty.
\]
Theorem~\ref{thm:main} gives the last claim, as $\dmax=d$.
\end{proof}

With $K=2$, $n=8m^{11}$ and $\eps\sim(n/8)^{-1/11}$, Corollary~\ref{cor:paths} already gives $L_1\ge(2-\eta)\dmax/\eps$. With $K_m=\lfloor\log m/(2\log12)\rfloor$ we have $12^{K}\le\sqrt m$, $\log n=(3K+5)\log m+\log 8$, and so $K\ge c\sqrt{\log n}$ for an absolute $c>0$ and large $m$. Since $H_n/E_n\to\infty$, this proves part~(b) of Theorem~\ref{thm:examples}; in this family $\eps=\exp(-(c'+o(1))\sqrt{\log n})$ with $c'=\sqrt{2\log 12/3}$.

The next family shows that the factor two persists however slowly $\eps$ tends to zero.

\begin{corollary}\label{cor:slow}
Let $s$ run through the odd integers and let $e_s\in(0,1)$ satisfy $e_s\to0$ and $e_s\ge s^{-10}$. Put $K=2$ and
\[
d=s^{20},\qquad M=s^{31},\qquad n=4s^{100},\qquad y=\lceil e_sn\rceil .
\]
Then the hypotheses of Theorem~\ref{thm:main} hold, $\eps=(8/3+o(1))\,e_s$, and $\PP(L_1\ge(2-\eta)\dmax/\eps)\to1$ for every $\eta\in(0,1)$.
\end{corollary}

\begin{proof}
Here $4\mid n$ and $4\mid(d-1)^2$ since $d$ is odd. As $e_sn\ge4s^{90}$, $y=(1+o(1))e_sn=o(n)$, and $(d-1)^2M\le s^{71}=o(n)$, so $N_1,N_3>0$ for large $s$, $\ell=(3/2+o(1))n$ and $\eps=4y/\ell=(8/3+o(1))e_s$.

\ref{C1}: $y=o(n)$ and $Md^3/n=s^{-9}/4\to0$. \ref{C2}: $n\eps^3\ge s^{70}\to\infty$ for large $s$, and $d\eps\ge s^{10}$ while $\log(n\eps^3)\le\log n=O(\log s)$. \ref{C3}: $n\eps\ge s^{90}$ for large $s$, while $K^2d^2=4s^{40}$ and $KMd^2=2s^{71}$. \ref{C4}: $\mu=M(M-1)d^2/(\ell-1)=(1+o(1))s^{102}/(6s^{100})\to\infty$. Apply Theorem~\ref{thm:main} with $K=2$ and $\eta/2$ in place of $\eta$.
\end{proof}

Part~(a) of Theorem~\ref{thm:examples} follows: given $\varphi$, apply Corollary~\ref{cor:slow} with $e_s=\max\{\varphi(4s^{100}),s^{-10}\}$ for large $s$. Then $e_s\to0$, and $\eps=(8/3+o(1))e_s\ge\varphi(n)$ for large $s$. Corollary~\ref{cor:slow} also allows $\eps$ to tend to zero as fast as a constant times $n^{-1/10}$.

\begin{remark}\label{rem:alln}
The families are defined along subsequences of $n$, which suffices to refute a statement about all degree sequences satisfying the hypotheses. If a sequence defined for every $n$ is wanted, one can use, at the remaining $n$, $N_3=\lfloor n/4\rfloor-\lceil e n\rceil$ vertices of degree three for a target $e=e(n)\to0$ with $en^{1/3}\to\infty$, and degree one otherwise, changing one vertex of degree one into degree two if the total degree is odd. Then $|\ell\eps-4\lceil en\rceil|\le4$, so $\eps\sim 8e/3$ and $n\eps^3\to\infty$, and the limit law and the three moments are as before.
\end{remark}

The exact identities of Lemma~\ref{lem:counts} and~\eqref{eq:driftC}, the integrality and the size of the error terms of both families at concrete parameters, and Lemma~\ref{lem:pairs} together with the formulas of Lemma~\ref{lem:moments} on all $945$ perfect matchings of a sequence of ten half-edges are checked by the script in~\cite{code}. These finite computations are not used in the proofs.

\section*{Acknowledgements and funding}
The research presented in this paper was supported by the European Research
Council (ERC) under the European Union's Horizon Europe research and
innovation programme (grant agreement No.~101041711), by the Simons
Foundation as part of the Collaboration on the Mathematical and Scientific
Foundations of Deep Learning, by Heights Labs, by Convex Nexus Capital, by
the Israel Science Foundation (grant number 2258/19), and by the Israel
Science Foundation (ISF Grant 4101/25).

\section*{Declaration of competing interest}
The author holds equity in Heights Labs and in Convex Nexus Capital and
has a management role at Convex Nexus Capital. These entities are included
in the funding declaration above.

\section*{Declaration of generative AI and AI-assisted technologies}
The paper comes out of Math Brain, an ongoing project on the use of AI systems for research mathematics. On 11 October 2026 an agent running in OpenAI's Codex (GPT-6.1 Sol) obtained the counterexample of Corollary~\ref{cor:paths} with $K_m=\max\{2,\lfloor\sqrt{\log m}\rfloor\}$ and its proof; a further Codex agent and a Claude agent (Anthropic) checked the proof. On 17 and 18 October 2026 Claude Opus~5.5, working in Claude Code, checked the proof line by line, stated it as Theorem~\ref{thm:main} for general parameters, added the family of Corollary~\ref{cor:slow}, and wrote the computer checks and the paper. The author directed the work and is responsible for the content of the paper and for decisions concerning its dissemination.

\section*{Data availability}
No data were used. The code is available at\\ \url{https://github.com/yspennstate/hub-paths-barely-subcritical}.

\begin{thebibliography}{99}

\bibitem{Azu67} K.~Azuma, Weighted sums of certain dependent random variables, \emph{T\^ohoku Math. J.} \textbf{19} (1967), 357--367.

\bibitem{CP23} M.~Coulson and G.~Perarnau, Largest component of subcritical random graphs with given degree sequence, \emph{Electron. J. Probab.} \textbf{28} (2023), 1--28, doi:10.1214/23-EJP921.

\bibitem{Jan08} S.~Janson, The largest component in a subcritical random graph with a power law degree distribution, \emph{Ann. Appl. Probab.} \textbf{18} (2008), 1651--1668.

\bibitem{Luc90} T.~{\L}uczak, Component behavior near the critical point of the random graph process, \emph{Random Structures Algorithms} \textbf{1} (1990), 287--310.

\bibitem{vdH17} R.~van der Hofstad, \emph{Stochastic Processes on Random Graphs}, lecture notes for the 47th Summer School in Probability, Saint-Flour, 2017, version of 12 July 2017, \url{https://www.win.tue.nl/~rhofstad/SaintFlour_SPoRG.pdf}.

\bibitem{vdH26} R.~van der Hofstad, \emph{Stochastic Processes on Random Graphs}, book draft of 4 August 2026, Section~8.8, \url{https://rhofstad.win.tue.nl/Saint_Flour_SPoRG_Grundlehre.pdf}.

\bibitem{code} Y.~Shmalo, \texttt{checks.py} and its output, \url{https://github.com/yspennstate/hub-paths-barely-subcritical}.

\end{thebibliography}
\end{document}
